\section{Other known cases}
\label{sec:other}

At last, we mention several other families where murmuration densities have been computed.

\subsection{Flying Hecke characters of imaginary quadratic fields}

Wang \cite{wang2025murmurations} computed murmuration density for (possibly trivial) Hecke characters of imaginary quadratic fields.
Let $\scF$ be the family of nontrivial Hecke characters of $\bQ(\sqrt{-D})$ for square-free $D > 3$, $D \equiv 3 \pmod{4}$.
For $\psi \in \scF$, let $a_\psi(p) = \sum_{\rN\frp = p} \psi(\frp)$ be the trace of Frobenius at $p$.
The completed $L$-function $\Lambda(s, \psi)$ satisfies the functional equation $\Lambda(1 - s, \psi) = \Lambda(s, \psi)$, so the root numbeer is always $+1$.

\begin{theorem}[Wang \cite{wang2025murmurations}]
    \label{thm:wang}
    Then the average of normalized trace $\sqrt{p} a_\psi(p)$ over $\psi \in \scF$ with $N_\psi = |D_\psi| \in [X, X + Y]$ is
    \begin{equation}
        \label{eqn:wangdensity}
        \bE_{\substack{\psi \in \scF \\ N_\psi \in [X, X + Y]}}[\sqrt{p} a_\psi(p)] = \frac{\sum_{\substack{\psi \in \scF \\ N_\psi \in [X, X + Y]}} \sqrt{p} a_\psi(p)}{\sum_{\substack{\psi \in \scF \\ N_\psi \in [X, X + Y]}} 1} = c(p) \sum_{1 \le m < 2 \sqrt{y}} \delta_m(p) M_m(y) + M_-(y) + \text{error}
    \end{equation}
    where
    \begin{align*}
        M_m(y) &= \frac{11\zeta(2)}{4A} \sqrt{\frac{y}{4y - m^2}} \vartheta(m), \quad \vartheta(m) = 2^{\omega(m) + \min\{v_2(m), 2\}} \prod_{\substack{q \mid y \\ q \text{ odd}}} \left(1 + \frac{2 q^2 + q - 1}{q^4 - 3q^2 - 2q + 2}\right)\\
        M_-(y) &= -\frac{11\pi}{12 A} \sqrt{y}\\
        c(p) &= \frac{p+1}{3p} \prod_{\ell > 2, \left(\frac{p}{\ell}\right) = 1} \left(1 - 2 \ell^{-2} - \frac{2 \ell^{-3}}{1 - \ell^{-2}}\right) \\
        \delta_m(p) &= \begin{cases}
            \mathds{1}_{\left(\frac{p}{q}\right) = 1} & m = q^k, \,q\text{ is odd prime} \\
            \1_{p \equiv 3 \Mod{4}} & m = 2 \\
            \1_{p \equiv 5 \Mod{4}} & m = 4 \\
            \1_{p \equiv 1 \Mod{8}} & m = 2^\nu, \,\nu \ge 3
        \end{cases}
    \end{align*}
\end{theorem}
See \cite[Theorem 1]{wang2025murmurations} for the missing definitions.
Note that the main term of \eqref{eqn:wangdensity} depends on the arithmetic of $p$.
However, the dependence on $p$ is explicit and $c(p) \delta_m(p)$ is almost periodic in $m$, where such an almost periodicity does not appear in other families.
In particular, he proved that taking local average over short intervals of $p$ removes the dependence on $p$.
Especially, taking local average in an interval of length $H$ so that the primes are equidistributed in $[X, X + H]$, in the sense that
\[
\sum_{\substack{X \le n \le X + H \\ n \equiv a \Mod{q}}} \Lambda(n) = \frac{H}{\varphi(q)} (1 + o(1))
\]
where $\Lambda$ is the von Mangoldt function (for example, $H = X^{\frac{1}{2} + \epsilon}$ is enough under RH), then we get the following theorem.

\begin{theorem}[{Wang \cite[Theorem 2]{wang2025murmurations}}]
    \label{thm:wanglocal}
    Let $P \sim X^{1 + \delta_p}$ be a prime with fixed $\delta_p \ge 0$.
    We have
    \begin{equation}
        \bE_{p \in [P, P + H]}\left[\bE_{\substack{\psi \in \scF \\ N_\psi \in [X, X + Y]}}\left[\sqrt{p} a_\psi(p)\right]\right] = \overline{c} \sum_{1 \le m < 2 \sqrt{y}} \overline{M}_m(y) + M_-(y) + o(1)
    \end{equation}
    where
    \begin{align*}
        \overline{c} &= \frac{1}{3} \prod_{\ell > 2} \left(1 - \ell^{-2} - \frac{\ell^{-3}}{1 - \ell^{-2}}\right),\\
        \overline{M}_m(y) &= \frac{11\zeta(2)}{4A} \sqrt{\frac{y}{4y - m^2}} \kappa(m), \quad \kappa(m) = 2^{\1_{2 \mid m}} \prod_{\substack{q \mid m \\ q \text{ odd}}} \left(1 + \frac{q^2}{q^4 - 2q^2 - q + 1}\right)
    \end{align*}
    and $M_-(y)$ is the same as in Theorem \ref{thm:wang}.
\end{theorem}


The main ingredients of the proof are orthogonality of characters and summation of class numbers in short intervals with class number formula, similar as Zubrilina's approach.

\subsection{Flying modular forms (in weight direction)}


Recall that Zubrilina computed murmuration density for a \emph{fixed weight $k$ and varying level $N$}.
In \cite{bober2023murmurations}, Bober, Booker, M. Lee, and Lowry-Duda considered the opposite case, where they fix the level $N = 1$ and vary the weight $k$.
In this case, the considered family of Hecke newforms whose \emph{analytic conductor}
\[
\cN(k) := \left(\frac{\exp \psi(k/2)}{2 \pi}\right)^2 \approx \left(\frac{k-1}{4 \pi}\right)^2 + O(1)
\]
are in certain range, where $\psi(x) = \Gamma'(x)/\Gamma(x)$ is the digamma function.

\begin{theorem}[{Bober--Booker--Lee--Lowry-Duda \cite[Theorem 1.1]{bober2023murmurations}}]
    Fix $\epsilon \in (0, \frac{1}{12})$, $\delta \in \{0, 1\}$, and a compact interval $E \subset \bR_{> 0}$ with $|E| > 0$.
    Let $K, H > 0$ with $K^{\frac{5}{6} + \epsilon} < H < K^{1 - \epsilon}$, and let $N = \cN(K)$.
    Then as $K \to \infty$, we have

    \begin{equation}
        \frac{
            \sum_{\substack{p \text{ prime} \\ p / N \in E}}
            \log p
            \sum_{\substack{k \equiv 2 \delta \Mod{4} \\ |k - K| \le H}}
            \sum_{f \in H_k(1)} \lambda_f(p)
        }
        {
            \sum_{\substack{p \text{ prime} \\ p / N \in E}} 
            \log p 
            \sum_{\substack{k \equiv 2 \delta \Mod{4} \\ |k - K| \le H}} 
            \sum_{f \in H_k(1)} 1
        } = \frac{(-1)^\delta}{\sqrt{N}} \left(\frac{\nu(E)}{|E|} + o_{E, \epsilon}(1)\right)
    \end{equation}
    where
    \begin{align}
        \nu(E) &= \frac{1}{\zeta(2)} \sum_{\substack{a, q \in \bZ_{>0} \\ (a, q) = 1 \\ q^2 / a^2 \in E}} \frac{\mu(q)^2}{\varphi(q)^2 \sigma(q)} \left(\frac{a}{q}\right)^{-3} \\
        &= \frac{1}{2} \sum_{t \in \bZ} \prod_{p \nmid t} \frac{p^2 - p - 1}{p^2 - p} \int_E \cos\left(\frac{2 \pi t}{\sqrt{y}}\right) \dd y
    \end{align}
    where the summation $\sum^{*}$ indicates that the terms occuring at the endpoints of $E$ are halved. 
\end{theorem}

The main tool for the proof is the (original) Eichler--Selberg trace formula that does not include Atkin--Lehner operators (e.g. \cite[Theorem 2.1]{child2022twist}).
Then apply class number formula to replace class numbers with the special values of Dirichlet $L$-functions at $s = 1$, which can be estimated under GRH.

\subsection{Flying Maass forms}

Booker, Lee, Lowry-Duda, Seymour-Howell, and Zubrilina computed murmuration densities for weight 0 and level 1 Maass forms \cite{booker2024murmurations}.
They considered a family of Maass forms where the spectral parameter ($R$ with $\lambda = \frac{1}{4} + R^2$) goes ot $\infty$, which is equivalent to the \emph{analytic conductor} $\cN(R)$ going to $\infty$.

\begin{theorem}[{Booker--Lee--Lowry-Duda--Seymour-Howell--Zubrilina \cite[Theorem 1.1]{booker2024murmurations}}]
    Let $E \subset \bR_{>0}$ be a fixed compact interval with $|E| > 0$.
    Let $R, H > 0$ with $R^{\frac{5}{6} + \delta} < H < R^{1 - \delta}$ for some $\delta > 0$ and $N = \cN(R)$.
    Assumming GRH for $L$-functions of Dirichlet characters and Maass forms, as $R \to \infty$ we have 
    \begin{equation}
        \frac{\sum_{\substack{p\text{ prime} \\ p / N \in E}} \log p \sum_{|r(f) - R| \le H} \epsilon(f) a_{f}(p)}{\sum_{\substack{p\text{ prime} \\ p / N \in E}} \log p \sum_{|r(f) - R| \le H} 1} \to \frac{1}{\sqrt{N}|E|}  \sideset{}{^{*}}\sum_{\frac{q^2}{a^2} \in E}  \frac{\mu(q)^2}{\varphi(q)^2 \sigma(q)} \left(\frac{a}{q}\right)^{-3}
    \end{equation}
    where the summation $\sum^{*}$ indicates that the terms occuring at the endpoints of $E$ are halved. 
\end{theorem}

Proof uses an explicit Selberg trace formula due to Str\"ombergsson in his unpublished work \cite{strombergsson}, which requires an analytic test function and cannot be compactly supported, where GRH is needed to control the cutoff error term.
The remaining proof is similar to the weight aspect case of the modular forms \cite{bober2023murmurations}.


